% ===== Джерела до розділу 11. Ключі з префіксом ch11: =====
% Позиції ядра (bib-core.tex) тут НЕ дублюються.

% --- Правила та мови виявлення ---
\bibitem{ch11:sigma} Patzke T. Sigma: Generic Signature Format for SIEM Systems. SigmaHQ Project Documentation. URL: \url{https://github.com/SigmaHQ/sigma}.

\bibitem{ch11:suricata} Open Information Security Foundation (OISF). Suricata: Open Source Intrusion Detection and Prevention Engine --- User Guide. URL: \url{https://docs.suricata.io/}.

\bibitem{ch11:snort} Roesch M. Snort --- Lightweight Intrusion Detection for Networks. \emph{Proceedings of the 13th USENIX Conference on System Administration (LISA '99)}. USENIX Association, 1999. P.~229--238.

% --- Інженерія виявлення ---
\bibitem{ch11:detection-eng} Bianco D. The Pyramid of Pain. Enterprise Detection \& Response Blog, 2013. URL: \url{http://detect-respond.blogspot.com/2013/03/the-pyramid-of-pain.html}.

\bibitem{ch11:alerting} Zimmerman C. Ten Strategies of a World-Class Cybersecurity Operations Center. McLean: MITRE Corporation, 2014. 353~p.

\bibitem{ch11:detection-maturity} Stillions R. The DML Model (Detection Maturity Level). ryan.stillions blog, 2014.

% --- Безфайлові атаки та living-off-the-land ---
\bibitem{ch11:fileless} Kumar S., et al. An Emerging Threat: Fileless Malware --- A Survey and Research Challenges. \emph{Cybersecurity}. 2020. Vol.~3. Art.~1.

\bibitem{ch11:lolbas} LOLBAS Project. Living Off The Land Binaries, Scripts and Libraries. URL: \url{https://lolbas-project.github.io/}.

\bibitem{ch11:powershell} Bohannon D., Holmes L. Revoke-Obfuscation: PowerShell Obfuscation Detection Using Science. \emph{Black Hat USA}. Las Vegas, 2017.

% --- Українські джерела ---
\bibitem{ch11:detection-ua} Розроблення правил виявлення кіберзагроз для систем керування подіями безпеки (SIEM). \emph{Кібербезпека: освіта, наука, техніка}. 2024. URL: \url{https://csecurity.kubg.edu.ua/index.php/journal/article/view/745}.

\bibitem{ch11:ioc-ua} Формування та застосування індикаторів компрометації у процесах виявлення кіберінцидентів. \emph{Захист інформації}. 2023. Т.~25, №~1.
